Un électron entre horizontalement, avec une vitesse de \qty{\v{v0}e6}{m/s}, entre deux plaques parallèles horizontales de \qty{\v{L}}{cm} de longueur. Entre les plaques règne un champ électrique uniforme de \qty{\v{E}}{N/C} dirigé vers le haut ($m_e = \qty{9.109e-31}{kg}$, $e = \qty{1.602e-19}{C}$). On néglige la gravité.

\begin{enumerate}[label=\alph*)]
    \item Déterminez le module et le sens de l'accélération de l'électron entre les plaques.
    \item Combien de temps l'électron passe-t-il entre les plaques ?
    \item De combien l'électron est-il dévié verticalement à la sortie des plaques ?
\end{enumerate}
